% Irodov Part IV §4.4 original 201--212.
% Source PDF pp.191--192 printed pp.182--183; original answer PDF 333--334.
\emwaveproblem{201}{182}
空气平行板电容器两极板为半径 $R=6.0$cm 的圆盘，接频率 $\omega=1000\,\mathrm{s^{-1}}$ 的交变电压。求电容器内部磁场能峰值与电场能峰值之比。
\exsolution
忽略边缘场，极板间均匀电场
$E(t)=E_m\cos\omega t$。
半径 $r\le R$ 的环路上，由安培—麦克斯韦式
$2\pi rB_\varphi/\mu_0=
\varepsilon_0\pi r^2\dot E$，
所以磁场振幅
$B_{\varphi,m}=\mu_0\varepsilon_0\omega rE_m/2$。
若板距为 $d$，
\[
W_{B,m}=\int_0^R\frac{B_{\varphi,m}^2}{2\mu_0}
(2\pi rd)\,\mathrm dr
=\frac{\pi d\mu_0\varepsilon_0^2\omega^2E_m^2R^4}{16}.
\]
电场能峰值
$W_{E,m}=\frac12\varepsilon_0E_m^2\pi R^2d$。
故
\[
\boxed{\frac{W_{B,m}}{W_{E,m}}
=\frac{\varepsilon_0\mu_0\omega^2R^2}{8}
=\frac{\omega^2R^2}{8c^2}
\simeq5.0\times10^{-15}}.
\]
\excheck 两种峰值不同时刻出现，但可分别取各自时间最大值；比例与场振幅和板距无关。

\emwaveproblem{202}{182--183}
无限长螺线管圆形截面半径
$R=6.0$cm，绕组通角频率
$\omega=1000\,\mathrm{s^{-1}}$ 的正弦交流电。
求管内感应电场能峰值与磁场能峰值之比。
\exsolution
理想长螺线管内部均匀
$B_z(t)=B_m\cos\omega t$。
由法拉第定律对半径 $r<R$ 的圆环积分，
$2\pi rE_\varphi=-\pi r^2\dot B_z$，
故
$E_{\varphi,m}=\omega rB_m/2$。
单位长度的感应电场能峰值
\[
W'_{E,m}=\int_0^R\frac12\varepsilon_0E_{\varphi,m}^2(2\pi r)\,\mathrm dr
=\frac{\pi\varepsilon_0\omega^2B_m^2R^4}{16}.
\]
磁场能峰值（单位长度）
$W'_{B,m}=\pi R^2B_m^2/(2\mu_0)$。
因而
\[
\boxed{\frac{W'_{E,m}}{W'_{B,m}}
=\frac{\varepsilon_0\mu_0\omega^2R^2}{8}
\simeq5.0\times10^{-15}}.
\]
\excheck 与第201题呈电磁对偶：电容器电场为主、螺线管磁场为主，微弱的伴随场能比相同。

\emwaveproblem{203}{183}
空气平板圆盘电容器电极缓慢充电，忽略边缘场。证明：穿过电容器侧表面的坡印廷能流等于电容器电场储能的增长率。
\exsolution
板间电场 $E(t)$ 沿轴向 $z$，极板圆半径为 $R$、板距为 $d$。
由位移电流导致的环向磁场
$H_\varphi(r)=\varepsilon_0r\dot E/2$。
侧面 $r=R$ 处
$\bm S=\bm E\times\bm H$，
$\bm e_z\times\bm e_\varphi=-\bm e_r$，
即能流朝内，大小
$S_{\rm in}=EH_\varphi(R)=\varepsilon_0R E\dot E/2$。
穿过侧面的向内功率
\[
P_{\rm in}=2\pi Rd\,S_{\rm in}
=\varepsilon_0E\dot E\,\pi R^2d.
\]
电场储能
$W_E=\frac12\varepsilon_0E^2(\pi R^2d)$，
故
\[
\boxed{P_{\rm in}=\frac{\mathrm dW_E}{\mathrm dt}}.
\]
\excheck 充电时坡印廷矢量从空间侧面向内流入电容器，而不是沿导线穿越极板；严格瞬时场分布仍取准静态近似。

\emwaveproblem{204}{183}
一根圆截面直导体通直流 $I$。取一段导体，其电阻为 $R$，求穿过该段导体侧面的坡印廷能流。
\exsolution
设导体半径为 $a$、长度为 $\ell$，
轴向电场
$E_z=IR/\ell$。
在侧面磁场强度
$H_\varphi=I/(2\pi a)$。
坡印廷矢量
$\bm S=\bm E\times\bm H
=-E_zH_\varphi\bm e_r$，
朝向导体内部。
由侧面积 $2\pi a\ell$，
\[
\boxed{P_{\rm in}=2\pi a\ell E_zH_\varphi
=2\pi a\ell\frac{IR}{\ell}\frac{I}{2\pi a}
=I^2R}.
\]
\excheck 侧向流入导体的电磁能正好转化为焦耳热；结果与所选导体半径、长度的单独数值无关。

\emwaveproblem{205}{183}
非相对论质子经加速电压 $U$ 后形成圆柱形均匀离子束，其电流为 $I$。求束流外部距轴线 $r$ 处的坡印廷矢量模与方向。
\exsolution
质子束速度由动能给定
$v=\sqrt{2eU/m_p}$。
线电荷密度
$\lambda=I/v$。
束外径向电场
$E_r=\lambda/(2\pi\varepsilon_0r)
=I/(2\pi\varepsilon_0vr)$；
磁场强度
$H_\varphi=I/(2\pi r)$。
对沿 $+z$ 运动的正质子束，
$\bm E=E_r\bm e_r$、
$\bm H=H_\varphi\bm e_\varphi$，
所以
$\bm S=E_rH_\varphi\bm e_z$。
因此
\[
\boxed{\bm S=
\frac{I^2}{4\pi^2\varepsilon_0r^2}
\sqrt{\frac{m_p}{2eU}}\;\bm e_z}.
\]
\excheck $\bm e_r\times\bm e_\varphi=\bm e_z$；
坡印廷矢量沿束流前进方向，且束外按 $r^{-2}$ 衰减。

\emwaveproblem{206}{183}
长直螺线管的绕组电流以缓慢速率增大。证明：螺线管内部磁能的增长率等于从侧面流入的坡印廷能流。
\exsolution
设螺线管半径 $R$、长度 $\ell$，
内部近似均匀磁场 $B_z(t)$。
法拉第定律给侧面环向电场
$E_\varphi(R)=-R\dot B_z/2$，
磁场强度
$H_z=B_z/\mu_0$。
坡印廷矢量
$\bm S=\bm E\times\bm H
=E_\varphi H_z\bm e_r$，
当 $\dot B_z>0$ 时其径向分量为负，指向管内。
侧面流入功率
\[
P_{\rm in}
=2\pi R\ell\left(\frac{R\dot B_z}{2}\frac{B_z}{\mu_0}\right)
=\frac{\pi R^2\ell}{\mu_0}B_z\dot B_z.
\]
内部磁能
$W_B=B_z^2(\pi R^2\ell)/(2\mu_0)$，故
\[
\boxed{P_{\rm in}=\dot W_B}.
\]
\excheck 与第203题类似，能量从导体周围的空间流入充能区域；螺线管足够长且缓慢变化保证边缘及辐射修正小。

\emwaveproblem{207}{183}
图4.39表示双导线直流回路的一段，箭头给出两根导线内的电流方向，下导线电势高于上导线（$\varphi_2>\varphi_1$）。依据坡印廷能流判断电源在图的哪侧。
\exsolution
两导线之间的静电场从高电势导线2指向低电势导线1。
按图4.39，利用右手定则判断两导线电流合成的磁场在两线之间的方向。
取 $\bm E\times\bm H$，该能流沿两导线之间的空隙\emph{从左向右}输送能量。
因此提供电磁能的一端在上游：
\[
\boxed{\text{电源位于图4.39的左侧}}.
\]
\excheck 判断的是\emph{能流方向}而非沿单根导线的电子漂移方向；本题依赖原图给定的两个电流箭头与电势高低。

\emwaveproblem{208}{183}
理想长直同轴电缆以忽略的有功电阻向负载传输电压 $V$、电流 $I$。外同轴导电薄壳为回流导体。求穿过电缆任一横截面的电磁能流。
\exsolution
内芯与外壳间径向电场
$E_r=V/[r\ln(b/a)]$，
其中 $a,b$ 为内外电极半径。
环向磁场强度
$H_\varphi=I/(2\pi r)$。
坡印廷矢量沿电缆轴向，
\[
S_z=E_rH_\varphi
=\frac{VI}{2\pi r^2\ln(b/a)}.
\]
对环状电磁场截面积积分：
\[
\Phi=\int_a^bS_z\,2\pi r\,\mathrm dr
=\frac{VI}{\ln(b/a)}\int_a^b\frac{\mathrm dr}{r}.
\]
因此
\[
\boxed{\Phi=VI}.
\]
\excheck 电磁能主要经内芯与外导体之间的介质空间向负载传播，而非简单地“在铜导线内部流动”。

\emwaveproblem{209}{183}
理想长直同轴电缆向交流负载供能，
端口电压
$V(t)=V_0\cos\omega t$，
电流
$I(t)=I_0\cos(\omega t-\varphi)$。
求横截面平均坡印廷能流。
\exsolution
在准 TEM 传播模型中，
每一横截面的瞬时能流与端口电功率相等：
$\Phi(t)=V(t)I(t)$。
用恒等式
$\cos A\cos B=\tfrac12[\cos(A-B)+\cos(A+B)]$，
得到
\[
\Phi(t)=\tfrac12V_0I_0
[\cos\varphi+\cos(2\omega t-\varphi)].
\]
一完整周期平均，双频项为零：
\[
\boxed{\overline\Phi=
\frac12V_0I_0\cos\varphi
=V_{\rm rms}I_{\rm rms}\cos\varphi}.
\]
\excheck 无功电流分量导致场能往返交换，但平均功率只由功率因数 $\cos\varphi$ 决定。

\emwaveproblem{210}{183}
证明在两个介质的分界面上，坡印廷矢量法向分量在适当边界条件下连续。
\exsolution
坡印廷矢量 $\bm S=\bm E\times\bm H$。
在\emph{无自由面电流且边界不存在面层吸收或释放功率}的通常介质界面，
由法拉第定律与安培定律的边界条件，
切向分量满足
$\bm E_{1t}=\bm E_{2t}$、
$\bm H_{1t}=\bm H_{2t}$。
坡印廷矢量法向分量只由切向场交叉积决定：
\[
S_{1n}=\bm n\cdot(\bm E_{1t}\times\bm H_{1t}),
\quad
S_{2n}=\bm n\cdot(\bm E_{2t}\times\bm H_{2t}),
\]
所以
\[
\boxed{S_{1n}=S_{2n}}.
\]
\excheck 若界面承载自由面电流或具有表面耗散，切向 $\bm H$ 可跳变，界面能量平衡出现面功率项，题述连续性\emph{不再无条件成立}；本稿明确了原书省略的物理假设。

\emwaveproblem{211}{183}
证明在质心不受外力的非相对论封闭带电粒子系统中，若所有粒子具有相同比荷 $q_i/m_i=\kappa$，则没有电\emph{偶极}辐射。
\exsolution
系统电偶极矩
$\bm p=\sum_iq_i\bm r_i
=\kappa\sum_im_i\bm r_i
=\kappa M\bm R_{\rm cm}$，
其中 $M=\sum_i m_i$ 为总质量。
封闭系统质心匀速运动，因此
$\ddot{\bm R}_{\rm cm}=0$。
由经典非相对论电偶极辐射公式
\[
P_{\rm dip}=
\frac{|\ddot{\bm p}|^2}{6\pi\varepsilon_0c^3},
\]
得到
\[
\boxed{\ddot{\bm p}=\kappa M\ddot{\bm R}_{\rm cm}=0,
\qquad P_{\rm dip}=0}.
\]
\excheck 结论只是\emph{无电偶极辐射}；
并不排除更高多极矩辐射，不能误写成所有形式的电磁辐射都为零。

\emwaveproblem{212}{183}
电子沿一直线作谐振动，
振幅 $a=0.10\,\mathrm{nm}$，
角频率 $\omega=6.5\times10^{14}\,\mathrm{s^{-1}}$。
求一周期平均电磁辐射功率。
\exsolution
坐标 $x=a\cos\omega t$，
加速度
$\ddot x=-a\omega^2\cos\omega t$。
非相对论拉莫尔辐射瞬时功率
$P(t)=e^2\ddot x^2/(6\pi\varepsilon_0c^3)$，
周期平均平方加速度
$\langle\ddot x^2\rangle=\tfrac12a^2\omega^4$。
所以
\[
\boxed{\overline P=
\frac{e^2a^2\omega^4}{12\pi\varepsilon_0c^3}
\simeq5\times10^{-15}\,\mathrm W}.
\]
\excheck 振动加速度的\emph{均方值}而非峰值代入平均功率；使用 $a=10^{-10}$m。
